% Part: incompleteness
% Chapter: representability-in-q
% Section: introduction

\documentclass[../../../include/open-logic-section]{subfiles}

\begin{document}

\olfileid{inc}{req}{int}
\olsection{تعارف}

نامکملیت دے قضیے اوہناں نظرییاں اُتے لاگو ہوندے نیں جیہناں وچ
قابلِ حساب فنکشناں بارے بنیادی گلاں بیان تے ثابت کیتیاں جا سکدیاں
نیں۔ اسیں ایہو جیہا اک نہایت کمزور نظریہ دساں گے جیہنوں «$\Th{Q}$»
آکھدے نیں (یا کدی کدی رافائل رابنسن دے ناں اُتے «رابنسن دا $Q$»)۔
اسیں دساں گے کہ کسے فنکشن دا $\Th{Q}$ وچ \emph{قابلِ نمائندگی}
ہون دا کیہ مطلب اے، تے فیر ایہ گل ثابت کراں گے:
\begin{quote}
  کسے فنکشن دی $\Th{Q}$ وچ نمائندگی اُدوں تے صرف اُدوں ہو سکدی اے
جدوں اوہ قابلِ حساب ہووے۔
\end{quote}
ایس نال سانوں قابلِ حسابیت دا اک ہور نمونہ ملدا اے۔ ایہ وی
وکھاواں گے کہ سیٹ $\Setabs{!A}{\Th{Q} \Proves !A}$ فیصلہ پذیر نہیں۔
ایس لئی رُکّن دے مسئلے نوں ایس سیٹ دی رکنیت دے فیصلے وچ
گھٹاؤاں گے۔ آخر تک اسیں ایس توں کِتے ودھ طاقتور نتیجے ثابت
کر چُکے ہواں گے۔

$\Th{Q}$ دی زبان حساب دی زبان اے۔ $\Th{Q}$ دے مسلّمے ہِٹھ دِتّے نیں۔
اوہناں نوں !!{identity} والی پہلے درجے دی منطق دے ہور مسلّمیاں
تے قاعدیاں نال رلا کے ورتنا اے۔
\begin{align*}
& \lforall[x][\lforall[y][(\eq[x'][y'] \lif \eq[x][y])]] \tag{$!Q_1$}\\
& \lforall[x][\eq/[\Obj 0][x']] \tag{$!Q_2$}\\
& \lforall[x][(\eq[x][\Obj 0] \lor \lexists[y][\eq[x][y']])] \tag{$!Q_3$}\\
& \lforall[x][\eq[(x + \Obj 0)][x]] \tag{$!Q_4$}\\
& \lforall[x][\lforall[y][\eq[(x + y')][(x + y)']]] \tag{$!Q_5$}\\
& \lforall[x][\eq[(x \times \Obj 0)][\Obj 0]] \tag{$!Q_6$}\\
& \lforall[x][\lforall[y][\eq[(x \times y')][((x \times y) + x)]]] \tag{$!Q_7$}\\
& \lforall[x][\lforall[y][(x < y \liff \lexists[z][\eq[(z' + x)][y]])]] \tag{$!Q_8$}
\end{align*}
ہر قدرتی عدد $n$ لئی عددی عبارت $\num{n}$ اوہ ٹرم
$\Obj{0}^{\prime\prime\ldots\prime}$ اے جیہدے وچ کل $n$ پرائم نشان
ہون۔ سو $\num{0}$ صرف
!!{constant}~$\Obj{0}$ اے، $\num{1}$،
$\Obj{0}'$ اے، $\num{2}$، $\Obj{0}''$ اے، تے ایویں اگّے۔

حساب دے نظریے وجوں $\Th{Q}$ \emph{بہت ای} کمزور اے۔ مثال لئی،
ایس وچ $\lforall[x][\eq/[x][x']]$ یا
$\lforall[x][\lforall[y][(x + y) = (y + x)]]$ ورگیاں سادیاں
گلاں وی ثابت نہیں ہو سکدیاں۔ اسیں ویکھاں گے کہ $\Th{Q}$ دے
دلچسپ ہون دی وڈی وجہ \emph{ایہو} کمزوری اے۔ اوہ بس اینا کُ
طاقتور اے کہ نامکملیت دا قضیہ لاگو ہو سکے۔ $\Th{Q}$ دے
دلچسپ ہون دی اک ہور وجہ ایہ اے کہ اوہدے مسلّمیاں دا سیٹ
\emph{متناہی} اے۔

$\Th{Q}$ توں طاقتور نظریہ، جیہنوں \emph{پیانو حساب} $\Th{PA}$
کہندے نیں،~$\Th{Q}$ وچ استقرا دا اک خاکہ شامل کرن نال ملدا اے۔
\[
(!A(\Obj 0) \land \lforall[x][(!A(x) \lif !A(x'))]) \lif \lforall[x][!A(x)]
\]
جتھے $!A(x)$ کوئی وی فارمولا اے۔ جے $!A(x)$ وچ $x$ توں علاوہ ہور !!{variable}s آزاد ہون،
تاں اسیں شروع وچ کلی مسوّر لا کے اوہناں ساریاں نوں پابند کر
دیندے آں تاں جے استقرا دے خاکے دی ایہ صورت !!a{sentence} ہووے۔
مثال لئی، جے $!A(x, y)$ وچ !!{variable}~$y$ وی آزاد
ہووے، تاں متعلقہ مثال ایہ اے:
\[
\lforall[y][((!A(\Obj 0) \land \lforall[x][(!A(x) \lif !A(x'))]) \lif
  \lforall[x][!A(x)])]
\]
استقرا دے خاکے دیاں ایہو جہیاں صورتاں ورت کے~$\Th{PA}$ دے
مسلّمیاں توں، $\Th{Q}$ دے مسلّمیاں نالوں، کِتے ودھ کجھ ثابت
کیتا جا سکدا اے۔ حقیقت وچ، قدرتی عدداں بارے
ایہوجیہے «قدرتی» بیان لبھّن لئی کافی محنت پینی اے
جیہڑے پیانو حساب وچ ثابت نہ ہو سکّن!{}

\begin{defn}
\ollabel{defn:representable-fn}
  قدرتی عدداں توں قدرتی عدداں ول فنکشن $f(x_0,\ldots,x_k)$ نوں
{\em $\Th{Q}$ وچ قابلِ نمائندگی} آکھدے آں جے کوئی فارمولا
$!A_f(x_0,\dots,x_k,y)$ ایہو جیہا ہووے کہ جدوں وی
$f(n_0,\dots,n_k) = m$ ہووے، تاں $\Th{Q}$ ایہ دونوں گلاں ثابت کرے:
\begin{enumerate}
\item\ollabel{defn:rep:a} $!A_f(\num{n_0}, \dots, \num{n_k}, \num{m})$
\item\ollabel{defn:rep:b} $\lforall[y][(!A_f(\num{n_0}, \dots,
\num{n_k}, y) \lif \num{m} = y)]$۔
\end{enumerate}
\end{defn}

ایس تعریف نوں ہور طرح وی بیان کر سکدے آں۔ مثال لئی ایس
تعریف دے برابر ایہ شرط وی رکھی جا سکدی اے کہ $\Th{Q}$ ایہ ثابت کرے: $\lforall[y][(!A_f(\num{n_0}, \dots,
\num{n_k}, y) \liff \eq[y][\num{m}])]$۔

\begin{thm}
\ollabel{thm:representable-iff-comp}
کسے فنکشن دی $\Th{Q}$ وچ نمائندگی اُدوں تے صرف اُدوں ہو سکدی اے
جدوں اوہ قابلِ حساب ہووے۔
\end{thm}

ایس قضیے دیاں دوہاں سمتاں ثابت کرنیاں نیں۔ نحو نوں عددی روپ
دے چکے ہوئیے تاں قابلِ نمائندگی توں قابلِ حساب ہون والی سمت
کافی سِدھی اے۔ دوجی سمت لئی ودھ کم کرنا پینا اے۔ خیال ایہ اے کہ اسیں «قابلِ حساب» نوں
ٹھیک معنی دین لئی «عمومی بازگشتی» نوں چُناں گے تے وکھاواں گے کہ ہر عمومی بازگشتی فنکشن دی~$\Th{Q}$ وچ
نمائندگی ہو سکدی اے۔ یاد کرو، فنکشن عمومی بازگشتی
ہوندا اے جے اوہ $\Zero$، جانشین فنکشن~$\Succ$ تے
پروجیکشن فنکشناں~$\Proj{n}{i}$ توں، ترکیب، بنیادی
بازگشت تے باقاعدہ کمترین قدر دی تلاش ورت کے، تعریف کیتا جا سکے۔
سو ایہ دکھاون دا اک طریقہ کہ ہر عمومی بازگشتی
فنکشن دی~$\Th{Q}$ وچ نمائندگی ہو سکدی اے، ایہ اے
کہ بنیادی فنکشناں دی نمائندگی ہو سکدی اے تے جدوں
کجھ فنکشناں دی نمائندگی ہو سکدی ہووے، تاں انہاں
توں ترکیب، بنیادی بازگشت تے باقاعدہ کمترین قدر دی تلاش نال
تعریف کیتے فنکشناں دی وی ہو سکے۔ دوجے لفظاں وچ، اسیں ایہ وکھا سکدے آں کہ بنیادی فنکشن قابلِ
نمائندگی نیں تے ترکیب، بنیادی بازگشت تے باقاعدہ کمترین قدر
دی تلاش نال بنے فنکشن وی قابلِ نمائندگی فنکشناں دے طبقے
وچ رہندے نیں۔ ایہ یقینی بناندا اے کہ ہر عمومی بازگشتی
فنکشن دی نمائندگی ہو سکدی اے۔

ایہ وکھانا اوکھا اے کہ قابلِ نمائندگی فنکشناں توں بنیادی
بازگشت نال بنے فنکشن وی قابلِ نمائندگی نیں۔ ایس قدم توں
بچن لئی پہلاں وکھاواں گے کہ بنیادی بازگشت توں بغیر وی
کم چل سکدا اے۔ یعنی ہر عمومی بازگشتی فنکشن بنیادی فنکشناں
توں صرف ترکیب تے باقاعدہ کمترین قدر دی تلاش نال تعریف
کیتا جا سکدا اے۔ ایس لئی وکھاواں گے کہ اک خاص باقاعدہ کمترین
قدر دی تلاش بنیادی بازگشت دا کم کر سکدی اے۔ ایہ طریقہ
ورتن لئی کجھ ہور بنیادی فنکشن شامل کرن پین گے۔ اوہ جمع، ضرب
تے برابری دے تعلق دا امتیازی فنکشن~$\Char{=}$ نیں۔ فیر قضیہ
ایہ وکھا کے ثابت کر سکدے آں کہ \emph{ایہناں} سارے بنیادی
فنکشناں دی~$\Th{Q}$ وچ نمائندگی ہو سکدی اے تے ترکیب تے
باقاعدہ کمترین قدر دی تلاش نال قابلِ نمائندگی فنکشناں توں
بنے فنکشن وی قابلِ نمائندگی نیں۔

\end{document}
